% Part: proof-theory
% Chapter: natural-deduction
% Section: rules-proofs

\documentclass[../../../include/open-logic-section]{subfiles}

\begin{document}

\olfileid[hi]{pt}{nat}{rp}

\olsection{नियम और \usetoken{P}{proof}}

हम कुछ परिभाषाएँ देकर और कुछ शब्दावली नियत करके आरंभ करते हैं।

\Log{N1c} और~\Log{N1i} के उन दो नियमों पर विचार करें जिनमें $!A \lif !B$
आता है।
\[
\bottomAlignProof
\AxiomC{$\Discharge{!A}{x}$}
\shortDeduce
\DeduceC{$!B$}
\DischargeRule{\Intro{\lif}}{x}
\UnaryInfC{$!A \lif !B$}
\DisplayProof
\qquad
\bottomAlignProof
\AxiomC{$!A \lif !B$}
\AxiomC{$!A$}
\RightLabel{\Elim{\lif}}
\BinaryInfC{$!B$}
\DisplayProof
\]
\Elim{\lif} नियम पर्याप्त सरल है। रेखा के ऊपर के !!{formula}s
\emph{आधार-वाक्य} हैं। बाईं ओर का !!{formula}, \emph{प्रमुख}
!!{formula}~$!A \lif !B$ है। सभी !!{elimination} नियमों में ऐसा एक आधार-वाक्य
होता है, सदा सबसे बाईं ओर। इसे अनुमान का \emph{मुख्य} आधार-वाक्य कहते हैं।
यदि अन्य हों, जैसा इस मामले में है, तो उन्हें \emph{गौण} आधार-वाक्य कहते
हैं। !!{introduction} नियमों के आधार-वाक्य को भी मुख्य आधार-वाक्य कहते हैं।
रेखा के नीचे का !!{formula}, \emph{निष्कर्ष} है।

\Intro{\lif} नियम का केवल एक आधार-वाक्य है और यहाँ निष्कर्ष प्रमुख सूत्र~$!A
\lif !B$ है। सभी !!{introduction} नियम ऐसे ही हैं: निष्कर्ष प्रमुख सूत्र है;
उसका मुख्य संकारक वही है जिसे ``प्रस्तुत'' किया जा रहा है।

संकेतन $\Discharge{!A}{x}$ का निम्न अर्थ है। किसी !!a{proof} में
\Intro{\lif} अनुमान को आधार-वाक्य के रूप में किसी !!a{formula}~$!B$ पर लागू
किया जाता है।~$!B$ पर समाप्त होने वाले उप-!!{proof} में कई आरंभिक
!!{formula}s हैं। इन्हें \emph{मान्यताएँ} कहते हैं।~$!A$ रूप की एक या अधिक
मान्यताओं से~$!B$ का !!^a{proof}, ऐसा !!a{proof} है कि $!A$ से~$!B$ का
तात्पर्य है। जब हम \Intro{\lif} नियम लागू करते हैं, तो~$!A \lif !B$
निष्कर्षित करते हैं और निष्कर्ष अब मान्यता~$!A$ पर निर्भर नहीं रहता। इसे
बताने के लिए~\Intro{\lif} अनुमान सहित !!{proof} में~$!A$ रूप की मान्यताओं को
\emph{मुक्त} या \emph{बंद} किया जाता है। नियम बताता है कि कौन-सी मान्यताएँ
मुक्त की जा सकती हैं। निष्कर्ष $!A \lif !B$ वाले \Intro{\lif} के मामले में
वे~$!A$ रूप की मान्यताएँ हैं, अर्थात् सशर्त के पूर्ववर्ती से मेल खाती हैं।
कभी-कभी यह बताना महत्त्वपूर्ण होता है कि किस नियम पर कौन-सी मान्यताएँ मुक्त
की गईं, और यह मुक्तिकरण चिह्न~$x$ से किया जाता है। महत्त्वपूर्ण है कि
मान्यताएँ मुक्त करने वाले नियमों को अपेक्षित रूप की \emph{सभी} मान्यताएँ,
बल्कि वास्तव में कोई भी मान्यता, मुक्त करना \emph{आवश्यक नहीं}। उदाहरणतः
निम्न सही है, यद्यपि पहला \Intro{\lif} कोई मान्यता मुक्त नहीं करता:
\[
\AxiomC{$\Discharge{!A}{y}$}
\DischargeRule{\Intro{\lif}}{x}
\UnaryInfC{$!B \lif !A$}
\DischargeRule{\Intro{\lif}}{y}
\UnaryInfC{$!A \lif (!B \lif !A)$}
\DisplayProof
\]
जब कोई नियम किसी मान्यता को मुक्त नहीं करता, तो हम मुक्तिकरण चिह्न छोड़ देते
हैं (जैसा यहाँ किया)।

!!{proof} में
\[
\AxiomC{$\Discharge{!A}{x}$}
\AxiomC{$\Discharge{!A}{y}$}
\RightLabel{\Intro{\land}}
\BinaryInfC{$!A \land !A$}
\RightLabel{\Elim{\land}}
\UnaryInfC{$!A$}
\DischargeRule{\Intro{\lif}}{x}
\UnaryInfC{$!A \lif !A$}
\DischargeRule{\Intro{\lif}}{y}
\UnaryInfC{$!A \lif (!A \lif !A)$}
\DisplayProof
\]
पहला \Intro{\lif} अनुमान बाईं मान्यता $!A^x$ को और दूसरा दाईं
मान्यता~$!A^y$ को मुक्त करता है। अर्थात्~$x$ से चिह्नित सभी मान्यताएँ,
चिह्न~$x$ वाले \Intro{\lif} द्वारा और~$y$ चिह्न वाली सभी मान्यताएँ, चिह्न
~$y$ वाले द्वारा मुक्त की जाती हैं। इसके काम करने के लिए यह महत्त्वपूर्ण है
कि समान चिह्न वाली सभी मान्यताएँ समान !!{formula} भी हों।

परिमाणक नियम कुछ अधिक जटिल हैं। उदाहरणतः \Log{N1i} में $\lexists$ के नियम
हैं:
\[
\bottomAlignProof
\AxiomC{$!A(t)$}
\RightLabel{\Intro{\lexists}}
\UnaryInfC{$\lexists[x][!A(x)]$}
\DisplayProof
\qquad
\bottomAlignProof
\AxiomC{$\lexists[x][!A(x)]$} 
\AxiomC{$\Discharge{!A(c)}{x}$}
\shortDeduce
\DeduceC{$!B$}
\DischargeRule{\Elim{\lexists}}{x}
\BinaryInfC{$!B$} \DisplayProof 
\bottomAlignProof
\] \Intro{\lexists} नियम में आधार-वाक्य~$!A(t)$ है, जहाँ $t$ कोई बंद पद,
अर्थात् चरों से रहित पद, है। ऐसा अनुमान सही है---अर्थात् वह योजनाबद्ध नियम
\Intro{\lexists} से मेल खाता है---यदि कोई !!{formula}~$!A$, कोई चर~$x$ और
कोई बंद पद~$t$ ऐसा हो कि निष्कर्ष $\lexists[x][!A]$ और आधार-वाक्य
~$\Subst{!A}{t}{x}$ हो। \Elim{\lexists} हमें~$!A(c)$ रूप की मान्यताएँ मुक्त
करने देता है, अर्थात् प्रत्येक अनुमान के लिए कोई !!{constant}~$c$ है और
$\Subst{!A}{c}{x}$ रूप की मान्यताएँ मुक्त की जा सकती हैं। किसी एक अनुमान की
सभी मुक्त मान्यताओं में वही~$c$ प्रयुक्त होना चाहिए। इस अचर को
\Elim{\lexists} अनुमान का \emph{आइगेनचर} कहते हैं। अनुमान तभी सही है जब~$c$
मुक्त की गई~$!A(c)$ के अतिरिक्त किसी खुली मान्यता में न आए,~$!B$ में न आए
और~$!A$ में न आए। इसे \emph{आइगेनचर शर्त} कहते हैं।

\Log{N1c} और~\Log{N1i} में !!^{proof}, अनुक्रमों के परिमित वृक्ष हैं जिन्हें
अनुमान नियमों से अभिगृहीतों से आगमनात्मक रूप से उत्पन्न किया जाता है। प्रत्येक
पत्ती एक मान्यता है, संभवतः मुक्तिकरण चिह्न से चिह्नित, और प्रत्येक अन्य
!!{formula} अपने ठीक ऊपर के !!{formula}s से प्रणाली के किसी अनुमान नियम
द्वारा निष्पन्न होता है। यदि किसी अनुमान के ऊपर वृक्ष में मान्यता के रूप में
किसी !!{formula} की उपस्थिति को उसके ऊपर के किसी अनुमान ने मुक्त नहीं किया,
तो उसे (उस अनुमान पर) \emph{खुला} कहते हैं।

\begin{defn}\ollabel{defn:proof}
\Log{N1c} और~\Log{N1i} में \emph{!!^{proof}}, निम्न प्रकार आगमनात्मक रूप से
!!{formula}s के परिमित वृक्ष हैं:
\begin{enumerate}
  \item\ollabel{defn:prf:assumption} $x$ या $y$ जैसे मुक्तिकरण चिह्न से
  चिह्नित प्रत्येक !!{formula} अपने आप में !!a{proof} है। एक ही !!{formula}
  से बने !!a{proof} में वह !!{formula} खुली मान्यता गिना जाता है।
  \item\ollabel{defn:prf:one-prem} यदि $R$ ऐसा नियम है जिसके एक, दो या तीन
  आधार-वाक्य $!A_1$, $!A_2$, $!A_3$, निष्कर्ष~$!A$ हैं, और $\delta_1$,
  $\delta_2$, $\delta_3$ क्रमशः $!A_1$, $!A_2$, $!A_3$ पर समाप्त होने वाले
  !!{proof}s हैं, तो
  \[
  \AxiomC{}
  \RightLabel{$\delta_1$}
  \DeduceC{$!A_1$}
  \RightLabel{$R$}
  \UnaryInfC{$!A$}
  \DisplayProof
\qquad
  \AxiomC{}
  \RightLabel{$\delta_1$}
  \DeduceC{$!A_1$}
  \AxiomC{}
  \RightLabel{$\delta_2$}
  \DeduceC{$!A_2$}
  \RightLabel{$R$}
  \BinaryInfC{$!A$}
  \DisplayProof
\qquad
  \AxiomC{}
  \RightLabel{$\delta_1$}
  \DeduceC{$!A_1$}
  \AxiomC{}
  \RightLabel{$\delta_2$}
  \DeduceC{$!A_2$}
  \AxiomC{}
  \RightLabel{$\delta_3$}
  \DeduceC{$!A_2$}
  \RightLabel{$R$}
  \TrinaryInfC{$!A$}
  \DisplayProof
  \]
  क्रमशः !!a{proof} है, बशर्ते !!{proof}s में मान्यता-चिह्न संगत हों (दो
  भिन्न !!{formula}s का समान मुक्तिकरण चिह्न न हो) और अनुमान~$R$ के सही
  अनुप्रयोग हों।
  \item नए !!{proof} में सबसे ऊपर की !!{formula} उपस्थितियाँ, यदि वे
  $\delta_1$, $\delta_2$ या~$\delta_3$ की खुली मान्यताएँ हैं, तो नियम द्वारा
  मुक्त उपस्थितियों को छोड़कर !!{proof} की \emph{खुली मान्यताएँ} गिनी जाती
  हैं। यदि नियम~$x$ (या $x\, y$) से चिह्नित है, तो \Intro{\lif}
  और~\Intro{\lnot} के लिए ये~$\delta_1$ में~$x$ से चिह्नित सभी खुली
  मान्यताएँ हैं; \Elim{\lexists} के लिए~$\delta_2$ में~$x$ से चिह्नित सभी
  खुली मान्यताएँ; और \Elim{\lor} के लिए~$\delta_2$ में~$x$ से चिह्नित तथा
  ~ $\delta_3$ में~$y$ से चिह्नित सभी खुली मान्यताएँ।
\end{enumerate}
\end{defn}

किसी !!a{proof} के आगमनात्मक निर्माण में प्रयुक्त !!{proof}s को उसके
\emph{उप-!!{proof}s} कहते हैं। किसी अनुमान के आधार-वाक्यों पर समाप्त होने
वाले उप-!!{proof}s को उस अनुमान के \emph{निकटस्थ उप-!!{proof}s} कहते हैं।
\Log{N2c} और \Log{N2i} के नियम \olref{tab:N1} में दिए गए हैं।

\subfile{rules-N1}

\begin{defn}
किसी !!a{proof}~$\delta$ की \emph{!!{height}}~$\pheight{\delta}$ इस प्रकार
आगमनात्मक रूप से परिभाषित है।
\begin{enumerate}
  \item यदि $\delta$ किसी मान्यता से बना है, तो $\pheight{\delta} = 0$।
  \item यदि $\delta$ एक आधार-वाक्य वाले अनुमान पर समाप्त होता है जिसका
  निकटस्थ उप-!!{proof}~$\delta_1$ है, तो $\pheight{\delta} =
  \pheight{\delta_1} + 1$।
  \item यदि $\delta$ दो आधार-वाक्यों वाले अनुमान पर समाप्त होता है जिसके
  निकटस्थ उप-!!{proof}s~$\delta_1$ और~$\delta_2$ हैं, तो $\pheight{\delta} =
  \max(\pheight{\delta_1}, \pheight{\delta_2}) + 1$।
  \item यदि $\delta$, \Elim{\lor} पर समाप्त होता है जिसके निकटस्थ
  उप-!!{proof}s~$\delta_1$, $\delta_2$ और~$\delta_3$ हैं, तो
  $\pheight{\delta} = \max(\pheight{\delta_1}, \pheight{\delta_2},
  \pheight{\delta_3}) + 1$।
\end{enumerate}
यदि किसी अनुक्रम~$S$ का ऊँचाई $\le n$ वाला !!a{proof} है, तो हम $\Proves[n]
S$ लिखते हैं।
\end{defn}

\begin{ex}
\Log{N1i} में कोई भी !!{formula}, खुली मान्यता के रूप में स्वयं से
!!a{proof} गिना जाता है। अतः
\[
!A \land !B^x
\]
एक !!a{proof}~$\delta_1$ है। उसकी ऊँचाई~$0$ है। $\delta_1$ से हम
~\Elim{\land} के दो रूपों में से एक लागू करके दो !!{proof}s~$\delta_2$
और~$\delta_3$ प्राप्त कर सकते हैं
\[
\AxiomC{$!A \land !B^x$}
\RightLabel{\Elim{\land}[2]}
\UnaryInfC{$!B$}
\DisplayProof
\qquad
\AxiomC{$!A \land !B^x$}
\RightLabel{\Elim{\land}[1]}
\UnaryInfC{$!A$}
\DisplayProof
\]
~$\delta_2$ और~$\delta_3$ के अंतिम अनुमान का निकटस्थ उप-!!{proof} समान है:
$!A \land !B$। दोनों !!{proof}s में $!A\land !B$ की उपस्थिति खुली मान्यता
है। हमारे पास $\pheight{\delta_1} = \pheight{\delta_2} = 1$ है।

नियम \Intro{\land} को~$!B \land !A$ रूप के निष्कर्ष को उचित ठहराने के लिए
$!B$ और $!A$ रूप के दो आधार-वाक्य चाहिए। अतः निम्न एक !!a{proof}~$\delta_4$
है:
\[
\AxiomC{$!A \land !B^x$}
\RightLabel{\Elim{\land}[2]}
\UnaryInfC{$!B$}
\LeftSubproofLabel{$\delta_2$}
\AxiomC{$!A \land !B^x$}
\RightLabel{\Elim{\land}[1]}
\UnaryInfC{$!A$}
\RightSubproofLabel{$\delta_3$}
\RightLabel{\Intro{\land}}
\BinaryInfC{$!B \land !A$}
\DisplayProof
\]
उसकी ऊँचाई $\pheight{\delta_4} = \max(\pheight{\delta_2},
\pheight{\delta_3})+1 = \max(1,1) + 1 = 2$ है। उसमें मान्यताओं के रूप में
$!A \land !B$ की दो उपस्थितियाँ हैं और दोनों खुली हैं।

अंततः हम $\delta_5$ पाने के लिए \Intro{\lif} लागू कर सकते हैं:
\[
\AxiomC{$\Discharge{!A \land !B}{x}$}
\RightLabel{\Elim{\land}[2]}
\UnaryInfC{$!B$}
\LeftSubproofLabel{$\delta_2$}
\AxiomC{$\Discharge{!A \land !B}{x}$}
\RightLabel{\Elim{\land}[1]}
\UnaryInfC{$!A$}
\RightSubproofLabel{$\delta_3$}
\RightLabel{\Intro{\land}}
\BinaryInfC{$!B \land !A$}
\DischargeRule{\Intro{\lif}}{x}
\UnaryInfC{$(!A \land !B) \lif (!B \land !A)$}
\DisplayProof
\]
उसकी ऊँचाई $\pheight{\delta_4} + 1 = 3$ है। \Intro{\lif} अनुमान~$x$ से
चिह्नित $!A \land !B$ रूप की सभी मान्यता-उपस्थितियों को मुक्त करता है,
अर्थात् निकटस्थ उपप्रमाण की दोनों पहले खुली मान्यताओं को। चूँकि वे ही
$\delta_5$ की एकमात्र मान्यताएँ हैं, उसकी कोई खुली मान्यता नहीं है और इसलिए
वह अपने अंतिम-!!{formula} $(!A \land !B) \lif (!B \land !A)$ \emph{का}
!!a{proof} है।
\end{ex}

\end{document}
